\documentclass[submit]{../harvardml}
\usepackage{../common}

\course{CS1810-S26}
\assignment{Homework \#2}
\duedate{February 27, 2026 at 11:59 PM}

\usepackage[OT1]{fontenc}
\usepackage[colorlinks,citecolor=blue,urlcolor=blue]{hyperref}
\usepackage{graphicx}
\usepackage{float}
\usepackage{subcaption}
\usepackage{fullpage}
\usepackage{amsmath}
\usepackage{todonotes}
\usepackage{amssymb}
\usepackage{framed}
\usepackage{color}
\usepackage{soul}
\usepackage{todonotes}
\usepackage{listings}
\usepackage{enumitem}
\usepackage{bm}
\usepackage{bbm}


\newcommand{\B}{\text{B}}
\newcommand{\Beta}{\text{Beta}}

\usepackage[mmddyyyy,hhmmss]{datetime}

\definecolor{verbgray}{gray}{0.9}

\lstnewenvironment{csv}{%
  \lstset{backgroundcolor=\color{verbgray},
  frame=single,
  framerule=0pt,
  basicstyle=\ttfamily,
  columns=fullflexible}}{}
  

%%%%%%%%%%%%%%%%%%%%%%%%%%%%%%%%%%%%%%%%%%%
%% Solution environment
\usepackage{xcolor}
\newenvironment{solution}{
    \vspace{2mm}
    \color{blue}\noindent\textbf{Solution}:
}{}
%%%%%%%%%%%%%%%%%%%%%%%%%%%%%%%%%%%%%%%%%%%


\begin{document}

\begin{center}
  {\Large Classification and Bias-Variance Trade-offs}\\
\end{center}

\subsection*{Introduction}

This homework is about classification, bias-variance trade-offs, and
uncertainty quantification.

The datasets that we will be working with relate to astronomical observations and loan applicants
The first dataset, found at \verb|data/planet-obs.csv|,
contains information on whether a planet was observed (as a binary
variable) at given points in time. This will be used in Problem 1. The
second dataset, available at \verb|data/hr.csv|, details different
loan applicants and their measured debt to income ratio and credit score. You will
work with this data in Problem 3.

As a general note, for classification problems we imagine that we have
the input matrix $\boldX \in \reals^{N \times D}$ (or perhaps they
have been mapped to some basis $\bm{\Phi}$, without loss of
generality) with outputs now ``one-hot encoded."  This means that if
there are~$K$ output classes, rather than representing the output
label $y$ as an integer~${1,2,\ldots,K}$, we represent $\boldy$ as a
``one-hot" vector of length~$K$. A ``one-hot" vector is defined as
having every component equal to 0 except for a single component which
has value equal to 1.  For example, if there are $K = 7$ classes and a
particular data point belongs to class 3, then the target vector for
this data point would be~$\boldy = [0,0,1,0,0,0,0]$.  We will define
$C_1$ to be the one-hot vector for the 1st class, $C_2$ for the 2nd
class, etc.  Thus, in the previous example $\boldy = C_3$. If there
are $K$ total classes, then the set of possible labels is $\{C_1
  \ldots C_K \} = \{C_k\}_{k=1}^K$.  Throughout the assignment we will
assume that each label $\boldy \in \{C_k\}_{k=1}^K$ unless otherwise
specified. The most common exception is the case of binary
classification ($K = 2$), in which case labels are the typical
integers $y \in \{0, 1\}$.

\subsection*{Resources and Submission Instructions}

% We encourage you to read CS181 Textbook's Chapter 3 for more information on linear classification, gradient descent, and classification in the discriminative setting. Read Chapter 2.8 for more information on the trade-offs between bias and variance.

In problems 1 and 3, you may use \texttt{numpy} or \texttt{scipy}, but not \texttt{scipy.optimize} or \texttt{sklearn}. Example code is given in the provided notebook. In past years, students have reported numerical stability issues for problem 3 and sometimes problem 1; we encourage students to not worry too much about this, as we are mindful of this in grading, but we recommend trying to run your code on \textbf{Google Colab} if you encounter issues (since this has resolved the issue for many students in the past).

Please type your solutions after the corresponding problems using this \LaTeX\ template, and start each problem on a new page. Please submit the writeup PDF to the Gradescope assignment `HW2'. Remember to assign pages for each question. \textbf{You must include any plots in your writeup PDF.} Please submit your \LaTeX file and code files to the Gradescope assignment `HW2 - Supplemental.' The supplemental files will only be checked in special cases, e.g. honor code issues, etc. Your files should be named in the same way as we provide them in the repository, e.g. \texttt{hw2.pdf}, etc.


\newpage
%%%%%%%%%%%%%%%%%%%%%%%%%%%%%%%%%%%%%%%%%%%%%
% Problem 1
%%%%%%%%%%%%%%%%%%%%%%%%%%%%%%%%%%%%%%%%%%%%%

\begin{problem}[Exploring Bias-Variance and Uncertainty, 30pts]
In this problem, we will explore the bias and variance of a few
different model classes when it comes to logistic regression, use cross validation for model selection, and
investigate two sources of predictive uncertainty in a synthetic
(made-up) scenario.

We are using a powerful telescope in the northern hemisphere to gather
measurements of some planet of interest. At certain times however, our
telescope is unable to detect the planet due to its positioning around
its star.  The data in \verb|data/planet-obs.csv| records the
observation time in the ``Time" column and whether the planet was
detected in the ``Observed" column (with the value 1 representing that
it was observed).  These observations were taken over a dark, clear
week, which is representative of the region.  Since telescope time is
expensive, we would like to build a model to help us schedule and find
times when we are likely to detect the planet.

\begin{enumerate}
\item Before we begin fitting any models, let us first derive the mathematical explanation for the bias-variance tradeoff. Let $\mathcal{D}$ denote the training dataset, while $\mathbf{x}$ is a fixed input (not in $\mathcal{D}$). Suppose the labels are generated via $$y = f(\mathbf{x}) + \varepsilon$$ where $f(\mathbf{x})$ is the true regression function and $\varepsilon$ is a noise term with $\mathbb{E}[\varepsilon|\mathbf{x}] = 0$ and $\text{Var}(\varepsilon|\mathbf{x}) = \sigma^2$. We'll denote the predictor learned from $\mathcal{D}$ as $\hat f_\mathcal{D}(\mathbf{x})$. Then, we have the prediction error as $$\text{MSE} = \mathbb{E}_{\mathcal{D}, y|\mathbf{x}}[(y - \hat f_\mathcal{D}(\mathbf{x}))^2]$$ We will see that this error decomposes into a noise term, a bias term, and a variance term. \begin{enumerate}
    \item Show that the expression for MSE above decomposes into $$\mathbb{E}_{y|\mathbf{x}}[(y - f(\mathbf{x}))^2] + \mathbb{E}_\mathcal{D}[(f(\mathbf{x}) - \hat f_\mathcal{D}(\mathbf{x}))^2] + \text{\{third term\}}$$ Hint: Add and subtract $f(\mathbf{x})$.

    \item Show that the third term from our final expression in part (b) simplifies to 0. Hint: Use Adam's Law/Law of Total Expectation.

    \item Show that the first term from our final expression in part (b) represents the noise in our data, i.e. $\mathbb{E}_{y|\mathbf{x}}[(y - f(\mathbf{x}))^2] = \sigma^2$. 

    \item Define the average predictor $\bar f(\mathbf{x}) = \mathbb{E}_\mathcal{D}[\hat f_\mathcal{D}(\mathbf{x})]$. Now show that the the second term from our final expression in part (b) represents the squared bias plus variance, i.e. $\mathbb{E}_\mathcal{D}[(f(\mathbf{x}) - \hat f_\mathcal{D}(\mathbf{x}))^2] = (f(\mathbf{x}) - \bar f(\mathbf{x}))^2 + \mathbb{E}_\mathcal{D}[(\bar f(\mathbf{x}) - \hat f_\mathcal{D}(\mathbf{x}))^2]$. 
\end{enumerate}

  \item Now let's explore the use of different bases to model our data. First, randomly split the data into 10 mini-datasets of size $N = 30$. We will use three different basis functions to model the data. We implemented basis 1 as $[1,t]$, where t is the data to help you get started. Then, choose two basis functions, which you'll use to transform the input data in the regression. (Hint: Examples of such basis functions were given in Lecture 2.) \textbf{Please write your chosen basis functions in the written section of your submission for this part of the question. }
  
  For each of these bases, fit a logistic regression model using sigmoid($\boldw^\top \boldsymbol\phi(t)$) to each dataset by using gradient descent to
        minimize the negative log-likelihood.  This means you will be
        running gradient descent 10 times for each basis, once for each
        dataset.

        Use the given starting values of $\boldw$ and a learning rate of $\eta=0.001$, take 1,000 update
        steps for each gradient descent run, and make sure to average the
        gradient over the data points at each step. These parameters,
        while not perfect, will ensure your code runs reasonably quickly.
  
  \item After consulting with a domain expert, we find that the probability of observing the planet is periodic as the planet revolves around its star---we are more likely to observe the planet when it is in front of its star than when it is behind it. In fact, the expert determines that observation follows the generating process $y \sim \text{Bern}(f(t))$, where $f(t) = 0.4 \times \cos(1.1t + 1) + 0.5$ for $t \in [0, 6]$ and $y \in \{0,1\}$. Note that we, the modelers, do not usually see the true data distribution. Knowledge of the true $f(t)$ is only exposed in this problem to allow for verification of the true bias.

        Use the given code to plot the true process versus your learned models. Include your plots in your solution PDF.

        \textbf{In no more than 5 sentences}, explain how bias and variance are reflected in the 3 types of curves on the graphs.  How do the fits of the individual and mean prediction functions change?  Keeping in mind that none of the model classes match the true generating process exactly, discuss the extent to which each of the bases approximates the true process.


    \item If we were to increase the size of each dataset drawn from $N = 30$ to a larger number, how would the bias and variance change for each basis? Why might this be the case? You may experiment with generating your own data that follows the true process and plotting the results, but this is \textbf{not} necessary. \textbf{Your response should not be longer than 5 sentences}.

    \item Consider the test point $t = 0.1$. Using your models trained on basis $\boldsymbol\phi_3$, report the predicted probability of observation of the \textit{first} model (the model trained on the first 30 data points). How can we interpret this probability as a measure of uncertainty? Then, compute the variance of the classification probability over your 10 models at the same point $t = 0.1$. How does this measurement capture another source of uncertainty, and how does this differ from the uncertainty represented by the classification probability? Repeat this process (reporting the first model's classification probability and the variance over the 10 models) for the point $t = 3.2$.

          Compare the uncertainties and their sources at times $t=0.1$ and $t=3.2$.

    \item We now need to make some decisions about when to request time on
          the telescope.  The justifications of your decisions will be sent to
          your funding agency, which will determine whether you will be
          allocated funds to use the telescope for your project. You seek out a team that has used the alternative telescope
                  for observing this planet, and they provide you their observation
                  file \verb|data/planet-obs-alternate.csv|.
                  Compare the observations from your telescope to theirs.  What
                  seems to be happening?  What might be an appropriate model for
                  this? Your funding agency asks you to refit your models on these
                  new data.  Do you think this is a reasonable ask, and if so, how
                  will it help you make better decisions about when to request
                  viewing time?  If not, why do you think the additional modeling
                  will not help? You do \emph{not} need to do any modeling; answer this question in \textbf{no more than 10 lines}. We are looking for your reasoning; there may be
          more than one valid answer.

  
      \item In part 1, we have selected 3 possible bases to use as features when training a model. This gave us 3 possible models. In this part, we will choose the right learning algorithm by using cross validation. Take your dataset split and train 10 models for each basis, using 9 of the chunks for training and 1 of the chunks for validation. Compare the average error across the 10 folds for the 3 bases. Which basis performs the best? Why do you think it performs the best?    

      
  \end{enumerate}
\end{problem}

\newpage

\begin{solution}

\section{Problem 1.}

\subsection{Question 1.}

\textbf{(a).} Let $x$ be fixed and $D$ be the (random) training dataset. Define $\hat f_D(x)$ as the predictor learned from $D$.
\[
\mathrm{MSE} \;=\; \mathbb{E}_{D,y\mid x}\!\left[(y-\hat f_D(x))^2\right].
\]
Add and subtract $f(x)$:
\[
y-\hat f_D(x) \;=\; (y-f(x)) + (f(x)-\hat f_D(x)).
\]
Thus,
\begin{align*}
\mathrm{MSE}
&=\mathbb{E}_{D,y\mid x}\!\left[\big((y-f(x))+(f(x)-\hat f_D(x))\big)^2\right] \\
&=\mathbb{E}_{D,y\mid x}\!\left[(y-f(x))^2\right]
+\mathbb{E}_{D,y\mid x}\!\left[(f(x)-\hat f_D(x))^2\right]
+2\,\mathbb{E}_{D,y\mid x}\!\left[(y-f(x))(f(x)-\hat f_D(x))\right] \\
&=\mathbb{E}_{y\mid x}\!\left[(y-f(x))^2\right]
+\mathbb{E}_{D}\!\left[(f(x)-\hat f_D(x))^2\right]
+\underbrace{2\,\mathbb{E}_{D,y\mid x}\!\left[(y-f(x))(f(x)-\hat f_D(x))\right]}_{\text{third term}},
\end{align*}
where we used that $(y-f(x))^2$ depends only on $y\mid x$ and $(f(x)-\hat f_D(x))^2$ depends only on $D$.
\\\\
\textbf{(b). }Using $y=f(x)+\varepsilon$, we have $y-f(x)=\varepsilon$ and $\mathbb{E}[\varepsilon\mid x]=0$. Then by iterated expectation,
\begin{align*}
\mathbb{E}_{D,y\mid x}\!\left[(y-f(x))(f(x)-\hat f_D(x))\right]
&=\mathbb{E}_{D}\!\left[(f(x)-\hat f_D(x))\,\mathbb{E}_{y\mid x}[\varepsilon\mid x]\right] \\
&=\mathbb{E}_{D}\!\left[(f(x)-\hat f_D(x))\cdot 0\right] \;=\;0.
\end{align*}
Hence the third term is $0$.
\\\\
\textbf{(c).} Since $y-f(x)=\varepsilon$,
\[
\mathbb{E}_{y\mid x}\!\left[(y-f(x))^2\right]
=\mathbb{E}[\varepsilon^2\mid x]
=\mathrm{Var}(\varepsilon\mid x)+(\mathbb{E}[\varepsilon\mid x])^2
=\sigma^2+0=\sigma^2.
\]
\\\\
\textbf{(d). }Define $\bar f(x)=\mathbb{E}_D[\hat f_D(x)]$. Then
\[
f(x)-\hat f_D(x)=(f(x)-\bar f(x))+(\bar f(x)-\hat f_D(x)).
\]
Squaring and taking $\mathbb{E}_D$,
\begin{align*}
\mathbb{E}_D[(f(x)-\hat f_D(x))^2]
&=(f(x)-\bar f(x))^2 + \mathbb{E}_D[(\bar f(x)-\hat f_D(x))^2]
+2(f(x)-\bar f(x))\,\mathbb{E}_D[\bar f(x)-\hat f_D(x)] \\
&=(f(x)-\bar f(x))^2 + \mathbb{E}_D[(\bar f(x)-\hat f_D(x))^2],
\end{align*}
since $\mathbb{E}_D[\bar f(x)-\hat f_D(x)]=\bar f(x)-\mathbb{E}_D[\hat f_D(x)]=0$.

\subsection{Question 2.}

We use the provided linear basis
\[
\phi_1(t) = \begin{bmatrix} 1 \\ t \end{bmatrix}.
\]
We choose the following two additional basis functions:
\[
\phi_2(t) =
\begin{bmatrix}
1 \\
t \\
t^2 \\
t^3
\end{bmatrix},
\qquad
\phi_3(t) =
\begin{bmatrix}
1 \\
\sin(1.1t) \\
\cos(1.1t) \\
\sin(2.2t) \\
\cos(2.2t)
\end{bmatrix}.
\]
For each basis, we trained logistic regression models of the form
\[
p(y=1 \mid t) = \sigma\!\left(w^\top \phi(t)\right),
\]
where $\sigma(z)=\frac{1}{1+e^{-z}}$, using gradient descent to minimize the negative log-likelihood. We used learning rate $\eta=0.001$, 1000 update steps, and averaged the gradient across data points at each step. Ten models were trained for each basis (one per mini-dataset).

\subsection{Question 3.}
\begin{figure}[H]
    \centering
    \includegraphics[width=0.5\linewidth]{hw2/img_output/basis1.png}
\end{figure}
\begin{figure}[H]
    \centering
    \includegraphics[width=0.5\linewidth]{hw2/img_output/basis2.png}
\end{figure}
\begin{figure}[H]
    \centering
    \includegraphics[width=0.5\linewidth]{hw2/img_output/basis3.png}
\end{figure}
Across the plots, the bias at a point $t$ is reflected by the vertical gap between the true function and the mean learned predictor, the variance is reflected by the spread of the individual learned curves around the mean. The linear basis $\phi_1$ exhibits high bias since it cannot capture periodic behavior, though the individual curves are relatively similar. The polynomial basis $\phi_2$ reduces bias by allowing more flexibility, but the individual curves vary more across datasets, indicating higher variance. The sinusoidal basis $\phi_3$ most closely matches the periodic structure of the true function and therefore has the lowest bias among the three. However, none of the bases match the true generating process exactly, so some residual bias remains in all cases.

\subsection{Question 4.}

If the dataset size $N$ increases, the variance of all three models decreases because the learned parameters become more stable across different samples of data. The bias, however, is primarily determined by the expressiveness of the basis and therefore does not substantially change with larger $N$. For the linear basis $\phi_1$, the model will continue to exhibit high bias since it cannot represent periodic structure, although its variance will shrink. For the polynomial basis $\phi_2$, variance will decrease noticeably with larger $N$, while the approximation bias remains moderate due to model misspecification. For the sinusoidal basis $\phi_3$, both variance decreases and bias is already relatively low, since its functional form is well aligned with the true periodic generating process.

\subsection{Question 5.}

For $t=0.1$, the first model predicts
\[
\hat p_1(0.1)=0.5654,
\]
and the variance of the predicted probabilities across the 10 models is
\[
\mathrm{Var}_{10}(0.1)=0.001323.
\]
The classification probability measures aleatoric uncertainty, since it reflects the model's estimated Bernoulli probability of observing the planet at that time.
The variance across models measures epistemic uncertainty, capturing how sensitive predictions are to the particular training dataset.
\\\\
For $t=3.2$, the first model predicts
\[
\hat p_1(3.2)=0.5053,
\]
with variance
\[
\mathrm{Var}_{10}(3.2)=0.001350.
\]
At both time points the epistemic uncertainty is small, but the aleatoric uncertainty differs: when the predicted probability is close to $0.5$ (as at $t=3.2$), intrinsic classification uncertainty is higher than when the probability is farther from $0.5$.

\subsection{Question 6.}

The alternate telescope’s observations appear to follow a similar periodic pattern but may differ in amplitude, phase, or baseline detection probability, suggesting a systematic shift between instruments. This could indicate measurement bias or differing observational sensitivity, rather than a change in the underlying planetary process. An appropriate model would incorporate telescope-specific effects. Refitting the models on the combined data is reasonable if the alternate data provide additional coverage or reduce sampling noise. More data can reduce variance in parameter estimates and improve calibration of detection probabilities. However, if the telescopes have fundamentally different noise characteristics, simply pooling data without modeling instrument effects could introduce bias. Accounting for these differences explicitly would allow more accurate prediction of high-probability observation windows. This would help allocate telescope time more efficiently by targeting periods with consistently high predicted detection probability across instruments.

\subsection{Question 7.}

The 10-fold average classification errors are:
\[
\text{basis1: }0.4967,\qquad \text{basis2: }0.3000,\qquad \text{basis3: }0.2800.
\]
Hence \(\phi_3\) (basis3) performs best. This is because the true observation probability is periodic, and \(\phi_3\) includes sinusoidal features at the relevant frequency, yielding lower approximation bias than \(\phi_1\) and typically a better bias-variance tradeoff than the higher-degree polynomial basis \(\phi_2\).

\end{solution}


\newpage
%%%%%%%%%%%%%%%%%%%%%%%%%%%%%%%%%%%%%%%%%%%%%
% Problem 2
%%%%%%%%%%%%%%%%%%%%%%%%%%%%%%%%%%%%%%%%%%%%%

\begin{problem}[Maximum likelihood in classification, 15pts]

Consider now a generative $K$-class model.  We adopt class prior
$p(\boldy = C_k; \bpi) = \pi_k$ for all $k \in \{1, \ldots, K\}$
(where $\pi_k$ is a parameter of the prior).
Let  $p(\boldx|\boldy=C_k)$ denote
the class-conditional density of features $\boldx$ (in this
case for class $C_k$). Consider the data set $\mathcal{D} = \{(\boldx_i, \boldy_i)\}_{i=1}^n$ where as above $\boldy_i \in \{C_k\}_{k=1}^K$ is
encoded as a one-hot target vector and the data are independent.

\begin{enumerate}
  \item Write out the log-likelihood of the data set, $\ln p(\mathcal{D} ; \bpi)$.

  \item Since the prior forms a distribution, it has the constraint that
        $\sum_k\pi_k - 1 = 0$.  Using the hint on
        Lagrange multipliers below, give the
        expression for the maximum-likelihood estimator for the prior
        class-membership probabilities, i.e.
        $\hat \pi_k.$
        Make sure to write out the intermediary equation you need
        to solve to obtain this estimator. Briefly state why your final answer is intuitive.
\end{enumerate}

For the remaining questions, let the
class-conditional probabilities be Gaussian distributions with
the same covariance matrix
$$p(\boldx | \boldy = C_k) = \mathcal{N}(\boldx |  \bmu_k, \bSigma), \text{\ for\ }k \in \{1,\ldots, K\}$$
and different means $\bmu_k$ for each class.

\begin{enumerate}
  \item[3.] Derive the gradient of the log-likelihood with respect to vector $\bmu_k$.
    Write the expression in matrix form as a function of the variables defined
    throughout this exercise. Simplify as much as possible for full credit.
  \item[4.] Derive the maximum-likelihood estimator $\hat{\mu}_k$ for vector $\bmu_k$. Briefly state why your final answer is intuitive.
  \item[5.] Derive the gradient for the log-likelihood with respect to the
    covariance matrix $\bSigma$ (i.e., looking
    to find an MLE for the covariance).
    Since you are differentiating with respect to a
    \emph{matrix}, the resulting expression should be a matrix!
    %
  \item[6.] Derive the maximum likelihood estimator $\hat{\Sigma}$ of the covariance matrix.
\end{enumerate}

\paragraph{Hint: Lagrange Multipliers.} Lagrange Multipliers are a method for
optimizing a function $f$ with respect to an
equality constraint, i.e.
\[\min_{\boldx} f(\boldx)\ \text{s.t.}\ g(\boldx) = 0.\]

This can be turned into an unconstrained problem by introducing a
Lagrange multiplier $\lambda$ and constructing the Lagrangian function,
\[L(\boldx, \lambda) =  f(\boldx) + \lambda g(\boldx).\]

It can be shown that it is a necessary condition that the optimum
is a critical point of this new function. We can find this point by solving two equations:

\[\frac{\partial L(\boldx, \lambda)}{\partial  \boldx} = 0  \ \ \text{and}\  \  \frac{\partial L(\boldx, \lambda)}{\partial \lambda} = 0 \]


\paragraph{Cookbook formulas.} Here are some formulas you might want to consider
using to compute difficult gradients. You can use them  in the homework
without proof. If you are looking to hone your matrix calculus skills, try to
find different ways to prove these formulas yourself (will not be part of the
evaluation of this homework). In general, you can use any formula from the matrix cookbook,
as long as you cite it. We opt for the following common notation:
$\boldX^{-\top} := (\boldX^{\top})^{-1}$
\begin{align*}
   & \frac{\partial \bolda^\top \boldX^{-1} \boldb}{\partial \boldX} = - \boldX^{-\top} \bolda \boldb^\top \boldX^{-\top} \\
   & \frac{\partial \ln | \det (\boldX) |}{\partial \boldX} = \boldX^{-\top}
\end{align*}

These formulas are from the matrix cookbook linked here: \href{https://www.math.uwaterloo.ca/~hwolkowi/matrixcookbook.pdf}{The Matrix Cookbook (University of Waterloo)}.  
\end{problem}


\newpage

\begin{solution}

\section{Problem 2.}

\subsection{Question 1.}

Let $D=\{(x_i,y_i)\}_{i=1}^n$ with one-hot $y_i\in\{C_1,\dots,C_K\}$, and prior $p(y=C_k;\pi)=\pi_k$. The data likelihood under the class prior is
\[
p(D;\pi)=\prod_{i=1}^n p(y_i;\pi)
=\prod_{i=1}^n \prod_{k=1}^K \pi_k^{\,y_{ik}},
\]
where $y_{ik}\in\{0,1\}$ is the $k$th component of the one-hot vector $y_i$.
Thus the log-likelihood is
\[
\ln p(D;\pi)=\sum_{i=1}^n\sum_{k=1}^K y_{ik}\ln \pi_k.
\]
\subsection{Question 2.}

Maximize $\sum_{i,k} y_{ik}\ln\pi_k$ subject to $\sum_{k=1}^K \pi_k-1=0$.
Let
\[
\mathcal{L}(\pi,\lambda)=\sum_{i=1}^n\sum_{k=1}^K y_{ik}\ln \pi_k+\lambda\Big(\sum_{k=1}^K \pi_k-1\Big).
\]
Setting derivatives to zero:
\[
\frac{\partial \mathcal{L}}{\partial \pi_k}=\frac{\sum_{i=1}^n y_{ik}}{\pi_k}+\lambda=0
\quad\Longrightarrow\quad
\pi_k=-\frac{\sum_{i=1}^n y_{ik}}{\lambda}.
\]
Using the constraint $\sum_k \pi_k=1$ gives
\[
1=\sum_{k=1}^K \pi_k=-\frac{1}{\lambda}\sum_{k=1}^K\sum_{i=1}^n y_{ik}
=-\frac{1}{\lambda}\sum_{i=1}^n \underbrace{\sum_{k=1}^K y_{ik}}_{=1}
=-\frac{n}{\lambda}
\quad\Longrightarrow\quad \lambda=-n.
\]
Hence the MLE is
\[
\hat\pi_k=\frac{1}{n}\sum_{i=1}^n y_{ik}.
\]
This is intuitive because $\hat\pi_k$ is exactly the empirical fraction of training points in class $k$.

\subsection{Question 3.}

For $p(x\mid y=C_k)=\mathcal{N}(x\mid \mu_k,\Sigma)$ with shared $\Sigma$, the (complete-data) log-likelihood is
\[
\ell(\{\mu_k\},\Sigma)=\sum_{i=1}^n\sum_{k=1}^K y_{ik}\left[-\frac12\log|\Sigma|-\frac12(x_i-\mu_k)^\top\Sigma^{-1}(x_i-\mu_k)\right]+ \text{const}.
\]
Using $\frac{\partial}{\partial \mu_k}\big[(x_i-\mu_k)^\top\Sigma^{-1}(x_i-\mu_k)\big]
=-2\Sigma^{-1}(x_i-\mu_k)$, we get
\[
\nabla_{\mu_k}\,\ell
=\sum_{i=1}^n y_{ik}\,\Sigma^{-1}(x_i-\mu_k)
=\Sigma^{-1}\!\left(\sum_{i=1}^n y_{ik}x_i-\Big(\sum_{i=1}^n y_{ik}\Big)\mu_k\right).
\]
Equivalently, letting $X=[x_1^\top;\dots;x_n^\top]\in\mathbb{R}^{n\times D}$ and $y_k=(y_{1k},\dots,y_{nk})^\top\in\mathbb{R}^n$,
\[
\nabla_{\mu_k}\,\ell=\Sigma^{-1}\!\left(X^\top y_k-(\mathbf{1}^\top y_k)\mu_k\right).
\]
\subsection{Question 4.}

Set the gradient to zero:
\[
0=\sum_{i=1}^n y_{ik}\,\Sigma^{-1}(x_i-\mu_k)
\;\Longrightarrow\;
\sum_{i=1}^n y_{ik}x_i=\Big(\sum_{i=1}^n y_{ik}\Big)\mu_k
\;\Longrightarrow\;
\hat\mu_k=\frac{\sum_{i=1}^n y_{ik}x_i}{\sum_{i=1}^n y_{ik}}.
\]
This is the sample mean of the points assigned to class $k$.

\subsection{Question 5.}

Using the matrix cookbook identities and $\mathrm{d}\,(x^\top\Sigma^{-1}x)= -\mathrm{tr}(\Sigma^{-1}xx^\top\Sigma^{-1}\mathrm{d}\Sigma)$, we obtain
\[
\nabla_{\Sigma}\,\ell
=-\frac12\sum_{i=1}^n\sum_{k=1}^K y_{ik}\,\Sigma^{-1}
+\frac12\sum_{i=1}^n\sum_{k=1}^K y_{ik}\,\Sigma^{-1}(x_i-\mu_k)(x_i-\mu_k)^\top\Sigma^{-1}.
\]
Since $\sum_{k=1}^K y_{ik}=1$ for one-hot labels, this simplifies to
\[
\nabla_{\Sigma}\,\ell
=-\frac{n}{2}\Sigma^{-1}
+\frac12\,\Sigma^{-1}\!\left(\sum_{i=1}^n\sum_{k=1}^K y_{ik}(x_i-\mu_k)(x_i-\mu_k)^\top\right)\Sigma^{-1}.
\]
\subsection{Question 6.}

Set $\nabla_{\Sigma}\ell=0$ and multiply on the left and right by $\Sigma$:
\[
-nI+\left(\sum_{i=1}^n\sum_{k=1}^K y_{ik}(x_i-\mu_k)(x_i-\mu_k)^\top\right)\Sigma^{-1}=0
\;\Longrightarrow\;
\hat\Sigma=\frac{1}{n}\sum_{i=1}^n\sum_{k=1}^K y_{ik}(x_i-\hat\mu_k)(x_i-\hat\mu_k)^\top.
\]


\end{solution}

\newpage
%%%%%%%%%%%%%%%%%%%%%%%%%%%%%%%%%%%%%%%%%%%%%
% Problem 3
%%%%%%%%%%%%%%%%%%%%%%%%%%%%%%%%%%%%%%%%%%%%%

\begin{problem}[Classifying Loan Applicants, 30pts]
In this problem, you will code up three different classifiers to classify different types of loan applicants. The file \verb|data/hr.csv| contains data on debt to income ratio measured in tenths of a percent and credit score. The data can be plotted on these two axes:
\begin{center}
  \includegraphics[width=.5\textwidth]{img_input/credit.png}
\end{center}
We've further transformed the raw data on debt to income ratio and credit score so that the default feature vector that you will be working with is defined as such:
\[\bm{x} = \left[\text{debt\_income\_ratio} \cdot \frac{200}{7}-7.5, \frac{\text{credit\_score}-500}{140}+0.5\right]^\top\]

\noindent Please implement the following classifiers:


\begin{enumerate}[label=\alph*)]

  \item \textbf{A generative classifier with Gaussian class-conditional
          densities with a \textit{shared covariance} matrix} across all classes.
        Feel free to re-use your Problem 2 results.

  \item \textbf{Another generative classifier with Gaussian class-conditional densities , but now
          with a \textit{separate covariance} matrix} learned for each class. (Note:
        The staff implementation can switch between the two Gaussian generative classifiers with just a
        few lines of code.)

  \item \textbf{A multi-class logistic regression classifier} using the softmax activation function. In your implementation of gradient descent, \textbf{make sure to use L2 regularization} with regularization parameter $\lambda = 0.001$. Please also include a \emph{\textbf{bias term, but do not regularize it}}. Limit the number of iterations of gradient descent to 200,000, and set the learning rate to be $\eta = 0.001$.

  \item \textbf{Another multi-class logistic regression classifier} with the additional feature map:
  $$\phi(\bm x) = [\ln (x_1+10), x_2^2]^\top$$

  \item \textbf{A kNN classifier} in which you classify based on the $k = 1$ and $k = 5$ nearest neighbors and the following distance function: 
  \[\text{dist}(\boldsymbol{x}, \boldsymbol{x}') = (x_1 - x'_1)^2/9 + (x_2 - x'_2)^2\]
        where nearest neighbors are those with the smallest distances from a given point.

        Note 1: When there are more than two labels, no label may have the
        majority of neighbors.  Use the label that has the most votes among
        the neighbors as the choice of label.

        Note 2: The grid of points for which you are making predictions
        should be interpreted as our test space.  Thus, it is not necessary
        to make a test point that happens to be on top of a training point
        ignore itself when selecting neighbors.

\end{enumerate}

After implementing the above classifiers, complete the following exercises:
  \begin{enumerate}



      \item Plot the decision boundaries generated by each classifier for the dataset. Include them in your PDF.
            Identify the similarities and differences among the classifiers. What explains the differences---in particular, which aspects or properties of each model dictate the shape of its decision boundary?
    
      \item
    
            Consider a loan applicant with Debt to Income Ratio 0.32 and Credit Score 350. To which class does each classifier assign this applicant? Report the classification probabilities of this applicant for models (c) and (d).
            
            Interpret how each model makes its classification decision. What else should we, the modelers, be aware of when making predictions on a point “far” from our training data? \textbf{Your response should no be longer than 5 sentences.}

    \item
        Can you think of any ethical problem that might arise from using this classifier to make loan decisions? You may approach this from any angle you like. For instance, can you think of someone who might have a low credit score and high debt-to-income ratio that you believe should nonetheless be offered a loan? Are there other variables that should be accounted for to ensure fair decisions? Are credit scores and debt-to-income ratio good bases for loan decisions? More generally, is using a classifier trained on past decisions to determine loan eligibility problematic in any way?
    \end{enumerate}
\end{problem}

\newpage

\begin{solution}

\section{Problem 3.}

\subsection{Question 1.}
\begin{figure}[H]
    \centering
    \includegraphics[width=0.5\linewidth]{hw2/img_output/p3_gaussian_shared.png}
\end{figure}
\begin{figure}[H]
    \centering
    \includegraphics[width=0.5\linewidth]{hw2/img_output/p3_gaussian_separate.png}
\end{figure}
\begin{figure}[H]
    \centering
    \includegraphics[width=0.5\linewidth]{hw2/img_output/p3_softmax_vanilla.png}
\end{figure}
\begin{figure}[H]
    \centering
    \includegraphics[width=0.5\linewidth]{hw2/img_output/p3_softmax_basis.png}
\end{figure}
\begin{figure}[H]
    \centering
    \includegraphics[width=0.5\linewidth]{hw2/img_output/p3_knn_k1.png}
\end{figure}
\begin{figure}[H]
    \centering
    \includegraphics[width=0.5\linewidth]{hw2/img_output/p3_knn_k5.png}
\end{figure}
The shared-covariance Gaussian classifier produces approximately linear decision boundaries separating the three classes, reflecting the affine discriminant functions induced by a common covariance matrix. The separate-covariance Gaussian classifier produces curved boundaries, visible in the bending region around the guarantor class, due to class-specific covariance matrices introducing second-order terms. The vanilla softmax classifier also yields linear decision boundaries, while the additional feature map significantly alters the geometry, producing a large horizontal band that classifies most middle credit-score values as requiring a guarantor. The kNN classifiers generate highly data-dependent boundaries: with $k=1$ the regions are more jagged and closely follow individual data points, whereas $k=5$ smooths the boundary and produces more stable regions. These differences arise from structural assumptions: generative Gaussian models impose distributional assumptions, softmax regression enforces linear separability in feature space, and kNN makes no parametric assumption and bases classification purely on local neighborhood structure.

\subsection{Question 2.}

All classifiers assign the applicant $(0.32, 350)$ to the Automatically Rejected class, as this point lies in the low credit-score region classified as rejection in all decision boundary plots. For models (c) and (d), the softmax classifiers compute posterior probabilities via the softmax function and select the class with highest probability; here, the rejection class has the largest score. The Gaussian classifiers compare class-conditional likelihoods under their covariance assumptions, both favoring rejection at this location. The kNN classifiers assign the label based on the majority of nearby training points, which are rejection examples. When predicting far from the training data, we must be cautious about extrapolation, as models may produce confident but unreliable predictions outside the data-supported region.

\subsection{Question 3.}

Using this classifier for loan decisions raises concerns about fairness and bias, since credit score and debt-to-income ratio may reflect structural inequalities rather than true repayment ability. A person with a temporarily low credit score due to medical debt or lack of credit history might still be financially stable and deserving of a loan. Relying only on these two variables ignores other relevant factors such as employment stability, savings, collateral, or recent improvements in financial behavior. Moreover, training a classifier on past decisions can perpetuate historical biases if earlier lending practices were discriminatory. Thus, without careful auditing and inclusion of broader contextual information, such models risk reinforcing inequities rather than making purely risk-based decisions.

\end{solution}

\newpage

\begin{problem}[Gradient Descent and Regularization, 15pts]

\noindent In this problem, you will analyze gradient descent for linear regression and study how \emph{ridge regularization} modifies the optimization landscape. You will first derive gradient descent updates for ordinary least squares, then extend them to the ridge-regularized objective. Finally, you will implement these algorithms to empirically observe how regularization affects convergence and generalization.

\medskip

\noindent Consider a regression dataset $\mathcal D = \{(\mathbf x_i, y_i)\}_{i=1}^N$, where each input $\mathbf x_i \in \mathbb{R}^D$ and response $y_i \in \mathbb{R}$. Let $\mathbf X \in \mathbb{R}^{N \times D}$ denote the design matrix and let $\mathbf y \in \mathbb{R}^N$ denote the vector of responses.
\\\\
\textbf{Detail:} The predictors in this dataset are \emph{highly correlated}.
As a result, the ordinary least squares objective may be ill-conditioned,
which can lead to unstable gradient descent updates and high-variance
parameter estimates. In later parts of this problem, we will examine how
ridge regularization modifies the optimization landscape in this setting.
\\\\
\noindent Consider the ordinary least squares objective
\begin{equation*}
    \mathcal L(\mathbf w) = \frac{1}{2} \|\mathbf X\mathbf w - \mathbf y\|_2^2.    
\end{equation*}

\begin{enumerate}

\item Derive the gradient $\nabla_{\mathbf w} \mathcal L(\mathbf w)$ with respect to the parameter vector $\mathbf w$. Your final expression should be written in matrix/vector form and simplified as much as possible.

\item Using your result, write the full-batch gradient descent update rule for
$\mathbf w$ with step size $\eta > 0$.
\\\\
\noindent Now consider the ridge-regularized objective
\begin{equation*}
    \mathcal L_\text{ridge}(\mathbf w) = \frac{1}{2} \|\mathbf X\mathbf w - \mathbf y\|_2^2 + \frac{\lambda}{2} \|\mathbf w\|_2^2,
\end{equation*}
where $\lambda \ge 0$ is a regularization parameter. Important note: typically, we do not regularize the bias term, $w_0$, thus $\|\mathbf w\|_2^2$ should be replace with $\mathbf w'$ where $w'_0 = 0$. However, for the sake of simplicity, we will omit that detail for the rest of the problem.

\item Derive the gradient $\nabla_{\mathbf w}\mathcal L_\text{ridge}(\mathbf w)$. Clearly indicate how the regularization term modifies the gradient compared to the unregularized case.


\item Rewrite the update in a form that makes the effect of ridge regularization
explicit. Briefly explain how this form illustrates the shrinkage effect of the
regularization term on the parameters.

\item You will implement gradient descent for both the ordinary least squares and
ridge-regularized objectives derived above and empirically study the effect of
regularization. Specifically, you will

\begin{itemize}
    \item Implement gradient and loss functions for both regular and ridge regression
    \item Write update step functions for stochastic gradient descent, stochastic gradient descent + momentum, and Adam optimizer.
\end{itemize}

Try not to touch anything that is not in between ``\textit{sol}" and ``\textit{los}"!


\item Generate the contour plots for the trajectories of the three algorithms across both the regularized and non-regularized case. What do you notice about the speed of convergence as well as where the algorithms converge to? What do you notice about the geometry of the loss function?\\
\textit{Note: the graphing code will pause when the first descent converges, so we see that some will not converge. Also, you'll see a subspace for the optimum in the unregularized case since we have a non-full rank data matrix $\mathbf X$!}

\end{enumerate}
\end{problem}

\pagebreak

\begin{solution}

\section{Problem 4.}

\subsection{Question 1.}

Let \(r(\mathbf{w}) = X\mathbf{w}-\mathbf{y}\). Then
\[
\mathcal{L}(\mathbf{w})=\frac12\|r(\mathbf{w})\|_2^2=\frac12 r(\mathbf{w})^\top r(\mathbf{w}).
\]
Using \(\nabla_{\mathbf{w}}\big(\tfrac12\|X\mathbf{w}-\mathbf{y}\|_2^2\big)=X^\top(X\mathbf{w}-\mathbf{y})\), we obtain
\[
\nabla_{\mathbf{w}}\mathcal{L}(\mathbf{w}) = X^\top(X\mathbf{w}-\mathbf{y}) = X^\top X\,\mathbf{w} - X^\top \mathbf{y}.
\]
\subsection{Question 2.}

With \(\nabla_{\mathbf{w}}\mathcal{L}(\mathbf{w}) = X^\top(X\mathbf{w}-\mathbf{y})\), the full-batch gradient descent update with step size \(\eta>0\) is
\[
\mathbf{w}^{(t+1)}=\mathbf{w}^{(t)}-\eta\,X^\top\!\big(X\mathbf{w}^{(t)}-\mathbf{y}\big)
=\big(I-\eta X^\top X\big)\mathbf{w}^{(t)}+\eta X^\top \mathbf{y}.
\]

\subsection{Question 3.}

For
\[
\mathcal{L}_{\text{ridge}}(\mathbf{w})=\frac12\|X\mathbf{w}-\mathbf{y}\|_2^2+\frac{\lambda}{2}\|\mathbf{w}\|_2^2,
\]
we have
\[
\nabla_{\mathbf{w}}\left(\frac12\|X\mathbf{w}-\mathbf{y}\|_2^2\right)=X^\top(X\mathbf{w}-\mathbf{y}),
\qquad
\nabla_{\mathbf{w}}\left(\frac{\lambda}{2}\|\mathbf{w}\|_2^2\right)=\lambda \mathbf{w}.
\]
Therefore,
\[
\nabla_{\mathbf{w}}\mathcal{L}_{\text{ridge}}(\mathbf{w})
= X^\top(X\mathbf{w}-\mathbf{y})+\lambda \mathbf{w}
= (X^\top X+\lambda I)\mathbf{w}-X^\top\mathbf{y}.
\]
Compared to the unregularized gradient \(X^\top(X\mathbf{w}-\mathbf{y})\), ridge adds the shrinkage term \(+\lambda \mathbf{w}\).

\subsection{Question 4.}

Using 
\[
\nabla_{\mathbf{w}}\mathcal{L}_{\text{ridge}}(\mathbf{w})
= X^\top(X\mathbf{w}-\mathbf{y})+\lambda \mathbf{w},
\]
the full-batch gradient descent update becomes
\[
\mathbf{w}^{(t+1)}
= \mathbf{w}^{(t)} - \eta \left[X^\top\!\big(X\mathbf{w}^{(t)}-\mathbf{y}\big) + \lambda \mathbf{w}^{(t)}\right].
\]
Rewriting,
\[
\mathbf{w}^{(t+1)}
= (1-\eta\lambda)\mathbf{w}^{(t)} 
- \eta X^\top\!\big(X\mathbf{w}^{(t)}-\mathbf{y}\big).
\]
This form makes the shrinkage explicit: before the data-driven update is applied, the parameter vector is multiplied by $(1-\eta\lambda)$, which contracts the weights toward zero at each iteration. Thus, ridge regularization continuously shrinks the magnitude of the coefficients, reducing variance and improving stability in ill-conditioned settings.

\subsection{Question 5.}

\textbf{Updated Code:}
\begin{lstlisting}[language=Python]
# -----------------------------
# Losses + gradients
# -----------------------------
def loss_linear(X, y, w):
    """Linear regression loss: L(w) = (1/2n) * sum_i (x_i^T w - y_i)^2"""
    n = X.shape[0]
    r = X @ w - y
    return 0.5 * np.mean(r ** 2)

def grad_linear(X, y, w):
    """Gradient of the linear regression loss w.r.t. w."""
    n = X.shape[0]
    r = X @ w - y
    return (X.T @ r) / n

def mse(X, y, w):
    """Mean squared error (no 1/2 factor). Used for plotting."""
    r = X @ w - y
    return np.mean(r ** 2)

def loss_ridge(X, y, w, lam):
    """Ridge loss: linear loss + (lambda/2) * ||w||^2"""
    return loss_linear(X, y, w) + 0.5 * lam * np.sum(w ** 2)

def grad_ridge(X, y, w, lam):
    """Gradient of ridge loss: grad_linear + lambda * w"""
    return grad_linear(X, y, w) + lam * w

# -----------------------------
# Optimizers (stateful)
# -----------------------------
@dataclass
class MomState:
    v: np.ndarray

def sgd_step(w, g, lr, state=None):
    """SGD update: w_{t+1} = w_t - lr * g_t"""
    w_new = w - lr * g
    return w_new, None

def mom_step(w, g, lr, beta, state: MomState | None):
    """
    Momentum:
        v_{t+1} = beta * v_t + g_t
        w_{t+1} = w_t - lr * v_{t+1}
    """
    if state is None:
        state = MomState(v=np.zeros_like(w))
    v_new = beta * state.v + g
    w_new = w - lr * v_new
    return w_new, MomState(v=v_new)

@dataclass
class AdamState:
    m: np.ndarray
    v: np.ndarray
    t: int

def adam_step(w, g, lr, beta1, beta2, eps, state: AdamState | None):
    """
    Adam:
        m_t = beta1 m_{t-1} + (1-beta1) g_t
        v_t = beta2 v_{t-1} + (1-beta2) g_t^2
        w_{t+1} = w_t - lr * m_hat / (sqrt(v_hat) + eps)
    """
    if state is None:
        state = AdamState(m=np.zeros_like(w), v=np.zeros_like(w), t=0)

    t = state.t + 1
    m = beta1 * state.m + (1 - beta1) * g
    v = beta2 * state.v + (1 - beta2) * (g ** 2)

    m_hat = m / (1 - beta1 ** t)
    v_hat = v / (1 - beta2 ** t)

    w_new = w - lr * m_hat / (np.sqrt(v_hat) + eps)
    return w_new, AdamState(m=m, v=v, t=t)
\end{lstlisting}

\subsection{Question 6.}
\begin{figure}[H]
    \centering
    \includegraphics[width=0.5\linewidth]{hw2/img_output/output.png}
\end{figure}
In the unregularized case (top row), the contours form a long flat valley due to the rank-deficient matrix $X$, so there is a subspace of minimizers along the dashed red line. SGD moves slowly and zig-zags along the narrow direction, while Momentum and Adam progress faster along the valley but may converge to different points on the same flat minimum set. In the ridge-regularized case (bottom row), the contours become elliptical with a unique minimum, eliminating the flat direction. All three algorithms now converge to the same point, and the trajectories are more stable. Thus, ridge regularization improves conditioning, speeds up convergence, and changes the geometry from a flat valley to a well-defined convex bowl.



\end{solution}


%%%%%%%%%%%%%%%%%%%%%%%%%%%%%%%%%%%%%%%%%%%%%
% Name and Calibration
%%%%%%%%%%%%%%%%%%%%%%%%%%%%%%%%%%%%%%%%%%%%%

\pagebreak
\textbf{Name}: Johnny Zhao

\textbf{Collaborators and Resources}: 
\begin{itemize}
    \item  AI Assistance: ChatGPT 5.2 (used only for high-level conceptual clarification and review of mathematical ideas).
    \item Course Materials: COMPSCI 1810 Machine Learning lecture notes and slides, BST 220 Applied Regression Analysis, BST 221 Applied Data Structures and Algorithms lecture notes, BST 236 Computing I textbook.
\end{itemize}


\end{document}
